\documentclass[
	ngerman,
	ruledheaders=section,%Ebene bis zu der die Überschriften mit Linien abgetrennt werden, vgl. DEMO-TUDaPub
	class=report,% Basisdokumentenklasse. Wählt die Korrespondierende KOMA-Script Klasse
	thesis={type=master},% Dokumententyp Thesis
	accentcolor=2c, % Akzentfarbe FG rtm
	custommargins=false,
	marginpar=false,
	BCOR=10mm, % Bindekorrektur (etwas größere mittlere Ränder)
	parskip=half-, % Absatzkennzeichnung durch Abstand vgl. KOMA-Sript
	fontsize=11pt, % Basisschriftgröße laut Corporate Design ist mit 9pt häufig zu klein
  	twoside,
  	numbers=noendperiod,
  	fleqn,
  	captions=oneline, % Einzeilige Captions zentrieren
	captions=tableheading, % Für besseren Abstand bei TabellenÜBERschriften
  	IMRAD=false,
]{tudapub}


%%%%%%%%%%%%%%%%%%%
%Pakete einbinden
%%%%%%%%%%%%%%%%%%%
%%%%%%%%%%%%%%%%%%%
% Mathe-Umgebungen und Symbole
%%%%%%%%%%%%%%%%%%%
\usepackage{amsmath}
\usepackage{amssymb}

%%%%%%%%%%%%%%%%%%%
% TikZ für Blockschaltbilder und Skizzen
%%%%%%%%%%%%%%%%%%%
\usepackage{tikz}				% Zum Erzeugen von Bildern mit TikZ
% TikZstyles-Bibliothek f�r Blockschaltbilder

\usetikzlibrary{arrows}		
\usetikzlibrary{arrows.meta}
\usetikzlibrary{positioning}

\newcommand{\TikZscale}{1}
\newcommand{\mm}{*\TikZscale mm}
\tikzset{every picture/.style={node distance=4\mm, >=stealth'}}

%% Bl�cke
	% Rechteckige Bl�cke
	\tikzstyle{block}      = [draw, semithick, rectangle, minimum height=8\mm, minimum width=8\mm, inner sep=3pt]
	\tikzstyle{NLblock}    = [draw, semithick, rectangle, minimum height=8\mm, minimum width=8\mm, inner sep=3pt, double distance=1.2pt]
	\tikzstyle{PICblock}   = [draw, semithick, rectangle, minimum height=8\mm, minimum width=8\mm, inner sep=2pt]
	\tikzstyle{NLPICblock} = [draw, semithick, rectangle, minimum height=8\mm, minimum width=8\mm, inner sep=3pt, double distance=1.2pt]
	\tikzstyle{noblock}	   = [rectangle, inner sep=-0.6pt]

	% Dreieckige Bl�cke
	\tikzstyle{Rgain}			 = [draw, semithick, isosceles triangle, inner sep=1pt, minimum height=8\mm, isosceles triangle apex angle=60]
	\tikzstyle{Lgain}			 = [draw, semithick, isosceles triangle, inner sep=1pt, minimum height=8\mm, isosceles triangle apex angle=60, shape border rotate=180]
	\tikzstyle{Ugain}			 = [draw, semithick, isosceles triangle, inner sep=1pt, minimum height=8\mm, isosceles triangle apex angle=60, shape border rotate=90]
	\tikzstyle{Dgain}			 = [draw, semithick, isosceles triangle, inner sep=1pt, minimum height=8\mm, isosceles triangle apex angle=60, shape border rotate=-90]

	% Runde Bl�cke
	\tikzstyle{sum}   	   = [draw, semithick, circle, inner sep=1pt, minimum size=3\mm]
	\tikzstyle{branch}		 = [draw, circle, inner sep=0pt, minimum size=1\mm, fill=black]
	\tikzstyle{BRANCH}		 = [coordinate]


%% Verbindungselemente
	% Linien mit Pfeil
	\tikzstyle{to}  		= [->, thick]
	\tikzstyle{toNL}		= [->, thick, shorten >=0.9pt]
	\tikzstyle{NLto}		= [->, thick, shorten <=0.9pt]
	\tikzstyle{NLtoNL}	= [->, thick, shorten <=0.9pt, shorten >=0.9pt]
	
	\tikzstyle{TO}  		= [semithick, double distance=2pt, shorten >=2mm, decoration={markings,mark=at position 1 with {\arrow[semithick]{open triangle 60}}}, preaction={decorate},postaction={draw, line width=2pt, white, shorten >= 1.5mm}]
	
	\tikzstyle{TONL}  		= [semithick, double distance=2pt, shorten >=2.2mm, decoration={markings,mark=at position 1 with {\arrow[semithick]{open triangle 60}}, transform={xshift=-0.7pt}}, preaction={decorate},postaction={draw, line width=2pt, white, shorten >= 1.5mm}]
	
	\tikzstyle{NLTO}  		= [semithick, double distance=2pt, shorten <=0.9pt, shorten >=2mm, decoration={markings,mark=at position 1 with {\arrow[semithick]{open triangle 60}}}, preaction={decorate}, postaction={draw, line width=2pt, white, shorten >= 1.5mm}]
	
	\tikzstyle{NLTONL}  		= [semithick, double distance=2pt, shorten <=0.9pt, shorten >=2.2mm, decoration={markings,mark=at position 1 with {\arrow[semithick]{open triangle 60}}}, preaction={decorate},postaction={draw, line width=2pt, white, shorten >= 1.5mm}]
	
	\tikzstyle{innerWhite} = [semithick, white,line width=2pt, shorten >= 2mm, shorten <= 2mm]
	
	%\tikzstyle{TO}  		= [semithick, double distance=2pt, shorten >=2mm, decoration={markings,mark=at position 1 with {\arrow[semithick]{open triangle 60}}}, postaction={decorate}]
		%
	%\tikzstyle{TONL} = [semithick, double distance=2pt, shorten >=2.2mm, decoration={markings,mark=at position 1 with {\arrow[semithick]{open triangle 60}}, transform={xshift=-0.7pt}}, postaction={decorate}]
	%
	%\tikzstyle{NLTO} = [semithick, double distance=2pt, shorten <=0.79pt, shorten >=2mm, decoration={markings,mark=at position 1 with {\arrow[semithick]{open triangle 60}}}, postaction={decorate}]
	%
	%\tikzstyle{NLTONL} = [semithick, double distance=2pt, shorten <=0.79pt, shorten >=2.2mm, decoration={markings,mark=at position 1 with {\arrow[semithick]{open triangle 60}}, transform={xshift=-0.7pt}}, postaction={decorate}]
	

	% Linien ohne Pfeil 
	\tikzstyle{line}       = [thick]
	\tikzstyle{lineNL}     = [thick, shorten >= 0.6pt]
	\tikzstyle{NLline}     = [thick, shorten <= 0.6pt]
	\tikzstyle{NLlineNL}	 = [thick, shorten <= 0.6pt, shorten >= 0.6pt]
	
	% Doppellinien ohne Pfeil
	\tikzstyle{LINE}       = [semithick, double distance=2pt]
	\tikzstyle{LINENL}     = [semithick, double distance=2pt, shorten >= 0.9pt]
	\tikzstyle{NLLINE}     = [semithick, double distance=2pt, shorten <= 0.9pt]
	\tikzstyle{NLLINENL}	 = [semithick, double distance=2pt, shorten <= 0.9pt, shorten >= 0.9pt]
	
	\tikzstyle{neg}				 = [postaction={decorate,decoration={markings, mark=at position 1 with{\draw[-](-2pt,-3pt)--(-2pt,-7pt);}}}]
	
	\tikzstyle{labelabove} = [above, anchor=base, yshift=0.75ex]
	\tikzstyle{labelbelow} = [below, anchor=base, yshift=-2ex]
	\tikzstyle{labelright} = [right]
	\tikzstyle{labelleft}  = [left]
	
	
	% Pins und Labels
	\tikzstyle{every pin edge}	= [<-, thick]
	\tikzstyle{every pin}	 			= [pin distance=5\mm]
	\tikzstyle{every label}			= [font=\small]
	\tikzstyle{terminal}				= [coordinate]
	
	% extra Symbole f�r circuits-Library
	\tikzstyle{jack}	= [draw, circle, minimum size=1.5mm, inner sep=0]
 % rtm-eigene Bibliothek für Blockschaltbilder

%%%%%%%%%%%%%%%%%%%
% Plotten direkt in LaTeX
%%%%%%%%%%%%%%%%%%%
\usepackage{pgfplots}

%%%%%%%%%%%%%%%%%%%
% Bilder und Abbildungen
%%%%%%%%%%%%%%%%%%%
\usepackage{pdfpages}
\usepackage{graphicx}
\usepackage{caption}
\usepackage{subcaption}
\usepackage{float}
\usepackage{import}
\usepackage{color}
\usepackage{transparent}


%%%%%%%%%%%%%%%%%%%
%Sprachanpassung & Verbesserte Trennregeln
%%%%%%%%%%%%%%%%%%%
%\usepackage{acronym}
\usepackage[acronym]{glossaries}
\setacronymstyle{long-short}
%\usepackage{glossaries-extra}
\usepackage[english, main=ngerman]{babel}
\usepackage[autostyle]{csquotes}% Anführungszeichen vereinfacht
\usepackage{microtype}
\usepackage{etoolbox}


%%%%%%%%%%%%%%%%%%%
%Meine Makros einbinden
%%%%%%%%%%%%%%%%%%%
\newcommand{\eqp}{\ensuremath{\, \, \, .}}
\newcommand{\M}[1]{\textbf{#1}} %Matrix  M
\newcommand{\Mr}[1]{\textbf{#1}_\rho} %Matrix mit tiefgestelltem rho, M_rho
\newcommand{\Mtil}[1]{\tilde{\textbf{#1}}} %Matrix mit Tilde
\newcommand{\Mtilt}[2]{\tilde{\textbf{#1}}_{#2}} %Matrix mit Tilde mit tiefgestelltem arg2
\DeclareRobustCommand{\Mt}[2]{\textbf{#1}_{#2}} %Matrix M mit tiefgestelltem arg2


\DeclareRobustCommand{\w}[1]{\underline{#1}} %Vektor arg1 unterstrichen 
\DeclareRobustCommand{\wt}[2]{\underline{#1}_{#2}} %Vektor unterstrichen w mit tiefgestelltem arg2
\DeclareRobustCommand{\wr}[1]{\underline{#1}_{\rho}} %Vektor unterstrichen mit rho, w_rho

% Textbausteine
% =============
	% Produktnamen
	\newcommand*{\Matlab}{\textsc{Matlab}}
	\newcommand*{\Matlabreg}{\textsc{Matlab}\textsuperscript{\tiny \textregistered}}
	\newcommand*{\MatSim}{\textsc{Matlab/Simulink}}
	\newcommand*{\Simulink}{\textsc{Simulink}}
	\newcommand*{\Simulinkreg}{\textsc{Simulink}\textsuperscript{\tiny \textregistered}}

	% Das Makro |\name|\marg{person} formatiert einen Personennamen bspw. eines Erfinders oder Entdeckers gemäß |\name{Euler}| \arrow\ \name{Euler}.
	\newcommand*{\name}[1]{\textsc{#1}}
	\newcommand*{\GenderPl}[1]{{#1}\_innen}
	\newcommand*{\GendeSin}[1]{{#1}\_in}
% Makros für Einheiten, Exponenten
% ================================

	\newcommand*{\unit}[1]{\ensuremath{\mathrm{#1}}}
	
	% Wert mit Einheit (mit kleinem Leerzeichen dazwischen), aus Text- UND Math-Modus
	\newcommand*{\valunit}[2]{\ensuremath{#1\,\mrm{#2}}}
	
	
	% "°C", im Text- oder Mathe-Modus
	\newcommand*{\degC}{
		\ifmmode
		^\circ \mrm{C}%
		\else
		\textdegree C%
		\fi}
	
	\newcommand*{\degree}{
		\ifmmode
		^\circ%
		\else
		\textdegree%
		\fi}
	
	% Für Exponentenschreibweise ( Anwendung: 123\E{3} )
	\newcommand*{\E}[1]{\ensuremath{\cdot 10^{#1}}}
	
	\newcommand*{\eexp}[1]{\ensuremath{\mathrm{e}^{#1}}}
	\newcommand*{\iu}{\ensuremath{\mathrm{j}}}
	
	\newcommand*{\todots}{\ensuremath{,\,\hdots,\,}}

% Makros für Formeln
% ==================

	% Definition für Vektor und Matizen
	\newcommand*{\mat}[1]{{\ensuremath{\boldsymbol{\mathrm{#1}}}}}
	\newcommand*{\ma}[1]{{\ensuremath{\boldsymbol{\mathrm{#1}}}}}
	\newcommand*{\mas}[1]{\ensuremath{\boldsymbol{#1}}}
	\newcommand*{\ve}[1]{\ensuremath{\boldsymbol{#1}}}
	\newcommand*{\ves}[1]{\ensuremath{\boldsymbol{\mathrm{#1}}}}
	
	\newcommand*{\AP}{\ensuremath{\mathrm{AP}}}
	\newcommand*{\doti}{\ensuremath{(i)^\cdot}}
	
	\newcommand*{\inprod}[2]{\ensuremath{\langle #1,\,#2 \rangle}}
	
	\newcommand*{\ul}[1]{\underline{#1}}
	
	% gerades "d" (z.B. für Integral)
	\newcommand*{\ud}{\ensuremath{\mathrm{d}}}
	
	% normaler Text in Formeln
	\newcommand*{\tn}[1]{\textnormal{#1}}
	
	% nicht-kursive Schrift in Formeln
	\newcommand*{\mrm}[1]{\ensuremath{\mathrm{#1}}}
	
	% gerades "T" für Transponiert
	\newcommand*{\transp}{\ensuremath{\mathrm{T}}}
	
	% gerades "rg"
	\newcommand*{\rang}{\ensuremath{\operatorname{rg}}}
	
	% Für geklammerte Ausdrücke mit Index (Subscript)
	% (einmal mit kursiven Index, einmal mit geradem Index)
	\newcommand*{\grpsb}[2]{\ensuremath{\left(#1\right)_{#2}}}
	\newcommand*{\grprsb}[2]{\ensuremath{\left(#1\right)_{\mathrm{#2}}}}
	
	% Ableitungen und Integrale
		% "normale" Ableitung (mit geraden "d"s)
		\newcommand*{\normd}[2]{\ensuremath{\frac{\mathrm{d}#1}{\mathrm{d}#2}}}
		\newcommand*{\normdat}[3]{\ensuremath{\left.\frac{\mathrm{d} #1}{\mathrm{d} #2}\right|_{#3}}}
		
		% Materielle Ableitung
		\newcommand*{\matd}[2]{\ensuremath{\frac{\mathrm{D} #1}{\mathrm{D} #2}}}
		\newcommand*{\matdat}[3]{\ensuremath{\left.\frac{\mathrm{D} #1}{\mathrm{D} #2}\right|_{#3}}}
		
		% Partielle Ableitung
		\newcommand*{\partiald}[2]{\ensuremath{\frac{\partial #1}{\partial #2}}}
		\newcommand*{\partialdat}[3]{\ensuremath{\left.\frac{\partial #1}{\partial #2}\right|_{#3}}}
	
	
	% Transformationen
	\newcommand*{\FT}[1]{\ensuremath{\mathfrak{F}\left\{#1\right\}}}
	\newcommand*{\FTabs}[1]{\ensuremath{\left|\mathfrak{F}\left\{#1\right\}\right|}}
	\newcommand*{\IFT}[1]{\ensuremath{\mathfrak{F}^{-1}\left\{#1\right\}}}
	\newcommand*{\DFT}[1]{\ensuremath{\mathrm{DFT}\left\{#1\right\}}}
	\newcommand*{\DFTabs}[1]{\ensuremath{\left|\mathrm{DFT}\left\{#1\right\}\right|}}
	\newcommand*{\Laplace}[1]{\ensuremath{\mathfrak{L}\left(#1\right)}}
	\newcommand*{\InvLaplace}[1]{\ensuremath{\mathfrak{L^{-1}}\left(#1\right)}}
	\newcommand*{\invtrans}{\ensuremath{\quad\bullet\!\!-\!\!\!-\!\!\circ\quad}}
	\newcommand*{\trans}{\ensuremath{\quad\circ\!\!-\!\!\!-\!\!\bullet\quad}}
	
	
	\newcommand*{\mlfct}[1]{\texttt{#1}}
	\newcommand*{\mlvar}[1]{\texttt{#1}}
	
	
	% Manche textcomp-Zeichen funktionieren mit dem TU-Design nicht, diese können dann mit diesem
	% Befehl gesetzt werden.
	\newcommand*{\textcompstdfont}[1]{{\fontfamily{cmr} \fontseries{m} \fontshape{n} \selectfont #1}}

% Makros für Variablen
% ================================
\newcommand*{\amax}{\ensuremath{a_{\textrm{max}}}}
\newcommand*{\amin}{\ensuremath{a_{\textrm{min}}}}
\newcommand*{\jmax}{\ensuremath{j_{\textrm{max}}}}
\newcommand*{\jmin}{\ensuremath{j_{\textrm{min}}}}
\newcommand*{\opt}[1]{\ensuremath{{#1}^*}}
\newcommand*{\fofx}{\ensuremath{f(\ve{x})}}
\newcommand*{\J}{\ensuremath{J(\ve{x}(t),\ve{u}(t),t)}}
\newcommand*{\Vofxoftf}{\ensuremath{V(\ve{x}(t_f),t_f)}}
\newcommand*{\Vofxoftone}{\ensuremath{\tilde{V}(\ve{x}(t_1),t_1)}}
\newcommand*{\Vofxoftfallgemein}{\ensuremath{\bar{V}(\ve{x}(t_f),t_f)}}
\newcommand*{\Vofxoftoftonefallgemein}{\ensuremath{\hat{V}(\ve{x}(t_f),\ve{x}(t_1),t_f,t_1)}}
\newcommand*{\variation}[1]{\ensuremath{\ve{\delta}_{#1}}}
\newcommand*{\dtint}[1]{\ensuremath{\textrm{d}{#1}}}
% Makros für x(t) Variablen
% ================================
\newcommand*{\xzero}{\ensuremath{\ve{x}_0}}
\newcommand*{\xoft}{\ensuremath{\ve{x}(t)}}
\newcommand*{\xoftf}{\ensuremath{\ve{x}(t_f)}}
\newcommand*{\xoftzero}{\ensuremath{\ve{x}(t_0)}}
\newcommand*{\xoftone}{\ensuremath{\ve{x}(t_1)}}
\newcommand*{\xofti}{\ensuremath{\ve{x}(t_i)}}
\newcommand*{\xoftoneminus}{\ensuremath{\ve{x}(t_1^-)}}
\newcommand*{\xoftoneplus}{\ensuremath{\ve{x}(t_1^+)}}
\newcommand*{\xoptoft}{\ensuremath{\opt{\ve{x}}(t)}}
\newcommand*{\xoptoftf}{\ensuremath{\opt{\ve{x}}(t_f)}}
\newcommand*{\xoptoftzero}{\ensuremath{\opt{\ve{x}}(t_0)}}
\newcommand*{\xoptoftfopt}{\ensuremath{\opt{\ve{x}}(\opt{t_f})}}
% Makros für x(tau) Variablen
% ================================
\newcommand*{\xoftau}{\ensuremath{\ve{x}(\tau)}}
\newcommand*{\xoftauf}{\ensuremath{\ve{x}(\tau_f)}}
\newcommand*{\xoftauzero}{\ensuremath{\ve{x}(\tau_0)}}
\newcommand*{\xoptoftau}{\ensuremath{\opt{\ve{x}}(t\tau)}}
\newcommand*{\xoptoftauf}{\ensuremath{\opt{\ve{x}}(\tau_f)}}
\newcommand*{\xoptoftauzero}{\ensuremath{\opt{\ve{x}}(\tau_0)}}
\newcommand*{\xoptoftaufopt}{\ensuremath{\opt{\ve{x}}(\opt{\tau_f})}}
% Makros für dx(t) Variablen
% ================================
\newcommand*{\dx}{\ensuremath{\dot{\ve{x}}}}
\newcommand*{\dxoft}{\ensuremath{\dot{\ve{x}}(t)}}
\newcommand*{\dxoftf}{\ensuremath{\dot{\ve{x}}(t_f)}}
\newcommand*{\dxoftzero}{\ensuremath{\dot{\ve{x}}(t_0)}}
\newcommand*{\dxoptoft}{\ensuremath{\dot{\ve{x}}^*(t)}}
\newcommand*{\dxoptoftf}{\ensuremath{\dot{\ve{x}}^*(t_f)}}
\newcommand*{\dxoptoftfopt}{\ensuremath{\dot{\ve{x}}^*(\opt{t_f})}}
% Makros für dx(tau) Variablen
% ================================
\newcommand*{\dxoftau}{\ensuremath{\ve{x}'(\tau)}}
\newcommand*{\dxoftauf}{\ensuremath{\ve{x}'(\tau_f)}}
\newcommand*{\dxoftauzero}{\ensuremath{\ve{x}'(\tau_0)}}
\newcommand*{\dxoptoftau}{\ensuremath{\ve{x}'^*(\tau)}}
\newcommand*{\dxoptoftauf}{\ensuremath{\ve{x}'^*(\tau_f)}}
\newcommand*{\dxoptoftaufopt}{\ensuremath{\ve{x}'^*(\opt{\tau_f})}}
% Makros für u(t) Variablen
% ================================
\newcommand*{\uoft}{\ensuremath{\ve{u}(t)}}
\newcommand*{\uoftf}{\ensuremath{\ve{u}(t_f)}}
\newcommand*{\uoftzero}{\ensuremath{\ve{u}(t_0)}}
\newcommand*{\uoptoft}{\ensuremath{\opt{\ve{u}}(t)}}
\newcommand*{\uoptoftf}{\ensuremath{\opt{\ve{u}}(t_f)}}
\newcommand*{\uoptoftfopt}{\ensuremath{\opt{\ve{u}}(\opt{t_f})}}
% Makros für u(tau) Variablen
% ================================
\newcommand*{\uoftau}{\ensuremath{\ve{u}(\tau)}}
\newcommand*{\uoftauf}{\ensuremath{\ve{u}(\tau_f)}}
\newcommand*{\uoftauzero}{\ensuremath{\ve{u}(\tau_0)}}
\newcommand*{\uoptoftau}{\ensuremath{\opt{\ve{u}}(\tau)}}
\newcommand*{\uoptoftauf}{\ensuremath{\opt{\ve{u}}(\tau_f)}}
\newcommand*{\uoptoftaufopt}{\ensuremath{\opt{\ve{u}}(\opt{\tau_f})}}
% Makros für lambda(t) Variablen
% ================================
\newcommand*{\lambdaoft}{\ensuremath{\ve{\lambda}(t)}}
\newcommand*{\lambdaoftf}{\ensuremath{\ve{\lambda}(t_f)}}
\newcommand*{\lambdaoftone}{\ensuremath{\ve{\lambda}(t_1)}}
\newcommand*{\lambdaoftzero}{\ensuremath{\ve{\lambda}(t_0)}}
\newcommand*{\lambdaoftoneminus}{\ensuremath{\ve{\lambda}(t_1^-)}}
\newcommand*{\lambdaoftoneplus}{\ensuremath{\ve{\lambda}(t_1^+)}}
% Makros für lambda(tau) Variablen
% ================================
\newcommand*{\lambdaoftau}{\ensuremath{\ve{\lambda}(\tau)}}
\newcommand*{\lambdaoftauf}{\ensuremath{\ve{\lambda}(\tau_f)}}
% Makros für dlambda(t) Variablen
% ================================
\newcommand*{\dlambda}{\ensuremath{\dot{\ve{\lambda}}}}
\newcommand*{\dlambdaoft}{\ensuremath{\dot{\ve{\lambda}}(t)}}
% Makros für dlambda(tau) Variablen
% ================================
\newcommand*{\dlambdaoftau}{\ensuremath{\ve{\lambda}'(\tau)}}
% Makros für z(t) Variablen
% ================================
\newcommand*{\zzero}{\ensuremath{\ve{z}_0}}
\newcommand*{\zoft}{\ensuremath{\ve{z}(t)}}
\newcommand*{\zoftf}{\ensuremath{\ve{z}(t_f)}}
\newcommand*{\zoftzero}{\ensuremath{\ve{z}(t_0)}}
\newcommand*{\zoptoft}{\ensuremath{\opt{\ve{z}}(t)}}
\newcommand*{\zoptoftf}{\ensuremath{\opt{\ve{z}}(t_f)}}
\newcommand*{\zoptoftzero}{\ensuremath{\opt{\ve{z}}(t_0)}}
\newcommand*{\zoptoftfopt}{\ensuremath{\opt{\ve{z}}(\opt{t_f})}}
% Makros für z(tau) Variablen
% ================================
\newcommand*{\zoftau}{\ensuremath{\ve{z}(\tau)}}
\newcommand*{\zoftauf}{\ensuremath{\ve{z}(\tau_f)}}
\newcommand*{\zoftauzero}{\ensuremath{\ve{z}(\tau_0)}}
\newcommand*{\zoptoftau}{\ensuremath{\opt{\ve{z}}(\tau)}}
\newcommand*{\zoptoftauf}{\ensuremath{\opt{\ve{z}}(\tau_f)}}
\newcommand*{\zoptoftauzero}{\ensuremath{\opt{\ve{z}}(\tau_0)}}
\newcommand*{\zoptoftaufopt}{\ensuremath{\opt{\ve{z}}(\opt{\tau_f})}}
% Makros für dz(t) Variablen
% ================================
\newcommand*{\dz}{\ensuremath{\dot{\ve{z}}}}
\newcommand*{\dzoft}{\ensuremath{\dot{\ve{z}}(t)}}
\newcommand*{\dzoftf}{\ensuremath{\dot{\ve{z}}(t_f)}}
\newcommand*{\dzoftzero}{\ensuremath{\dot{\ve{z}}(t_0)}}
\newcommand*{\dzoptoft}{\ensuremath{\dot{\ve{z}}^*(t)}}
\newcommand*{\dzoptoftf}{\ensuremath{\dot{\ve{z}}^*(t_f)}}
\newcommand*{\dzoptoftfopt}{\ensuremath{\dot{\ve{z}}^*(\opt{t_f})}}
% Makros für dz(tau) Variablen
% ================================
\newcommand*{\dzoftau}{\ensuremath{\ve{z}'(\tau)}}
\newcommand*{\dzoftauf}{\ensuremath{\ve{z}'(\tau_f)}}
\newcommand*{\dzoftauzero}{\ensuremath{\ve{z}'(\tau_0)}}
\newcommand*{\dzoptoftau}{\ensuremath{\ve{z}'^*(\tau)}}
\newcommand*{\dzoptoftauf}{\ensuremath{\ve{z}'^*(\tau_f)}}
\newcommand*{\dzoptoftaufopt}{\ensuremath{\ve{z}'^*(\opt{\tau_f})}}
% Makros für Ttransformation Variablen
% ================================
\newcommand*{\Ttransoftau}{\ensuremath{\mathcal{T}(\tau)}}
\newcommand*{\Ttransoftauzero}{\ensuremath{\mathcal{T}(\tau_0)}}
\newcommand*{\Ttransoftauf}{\ensuremath{\mathcal{T}(\tau_f)}}
% Makros für dTtransformation Variablen
% ================================
\newcommand*{\dTtransoftau}{\ensuremath{\mathcal{T}'(\tau)}}
\newcommand*{\dTtransoftauzero}{\ensuremath{\mathcal{T}'(\tau_0)}}
\newcommand*{\dTtransoftauf}{\ensuremath{\mathcal{T}'(\tau_f)}}



%%%%%%%%%%%%%%%%%%%
%Literaturverzeichnis
%%%%%%%%%%%%%%%%%%%
\catcode30=12  % Workaround for current biblatex problem (https://tex.stackexchange.com/questions/564990/error-after-miktex-reinstall-text-line-contains-an-invalid-character)
\usepackage[backend=biber,style=ieee,bibencoding=utf8]{biblatex}
\addbibresource{./Bib/literature.bib}

% ======================================================
% Fix für Fehlermeldungen 'Undifined control sequence' und 'Missing number, treated as zero.' bei gestapelten Akzenten in der Mathe-Umgebung. (z.B. \dot{\hat{x}})
% Der Fehler tritt auf bei tuda-ci v3.10 (Feb. 2021) in Kombination mit dem amsmath package.
%
% BITTE PRÜFEN, OB DER FEHLER NOCH BESTEHT UND DEN WORKAROUND GGF. WIEDER ENTFERNEN
%
% workaround/solution by skillmon posted here
% https://tex.stackexchange.com/questions/482281/math-accent-symbol-over-parentheses-enclosing-accented-symbol-amsmath
\makeatletter
\protected\def\mathaccentV#1#2#3#4#5%
  {%
    \ifmmode
      \mathaccentV@do{#2}{#3}{#4}{#5}%
    \else
      \@xp\nonmatherr@\csname #1\endcsname
    \fi
  }
\def\mathaccentV@do#1#2#3#4%
  {%
    \global\let\macc@nucleus\@empty
    \mathaccent"\accentclass@#1#2#3{#4}\macc@nucleus
  }
\makeatother
% ======================================================

\begin{document}
% =================================================================================
% Header mit allen Metadaten etc. 
% =================================================================================
\Metadata{	% Metadaten im erzeugten PDF/A
	title=Masterthesis Markus Amann,
	author=Markus Amann,
	keywords={TU Darmstadt \sep Corporate Design \sep LaTeX},
}

\title{Trajektorienplanung als dynamisches Optimalsteuerungsproblem und Interpretation des Lösungsraums für Fahrkomfort im automatisierten Fahren}
\subtitle{\textmd{Studiengang Elektrotechnik und Informationstechnik}}
\author[Markus Amann]{Markus Amann}
\reviewer*[Erstprüfer,Betreuer]{Prof. Dr.-Ing. Ulrich Konigorski \and Wi Mi M.Sc. Alexander Steinke} % erst ab Version v3.06 (2020-10-20) des tuda-ci-Pakets möglich

\department{etit}

\addTitleBox{\includegraphics[width=\linewidth]{Bilder/rtm_mit_schrift}}

%\date{\today}
\submissiondate{4. Juli 2022}

\lowertitleback{
	Technische Universität Darmstadt\\
	Institut für Automatisierungstechnik und Mechatronik\\
	Fachgebiet Regelungstechnik und Mechatronik\\
	Prof. Dr.-Ing. Ulrich Konigorski\\
}

\frontmatter
\maketitle

\clearpage
\setcounter{page}{1}
\section*{Aufgabenstellung}

%\documentclass[
%	paper=a4,
%	ngerman,
%	accentcolor=2c, % Akzentfarbe FG rtm
%	type=announcement,
%	marginpar=false,
%	fleqn,
%	class=report,
%	fontsize=11pt
%	]{tudaposter}
%
%%%%%%%%%%%%%%%%%%%%
%%Sprachanpassung & Verbesserte Trennregeln
%%%%%%%%%%%%%%%%%%%%
%\usepackage[english, main=ngerman]{babel}
%\usepackage[autostyle]{csquotes}% Anführungszeichen vereinfacht
%\usepackage{microtype}
%
%\usepackage{amsmath,amssymb}
%
%\usepackage{url}
%
%%%%%%%%%%%%%%%%%%%%
%%Literaturverzeichnis
%%%%%%%%%%%%%%%%%%%%
%\catcode30=12  % Workaround for current biblatex problem (https://tex.stackexchange.com/questions/564990/error-after-miktex-reinstall-text-line-contains-an-invalid-character)
%\usepackage[backend=biber,style=ieee,bibencoding=utf8]{biblatex}
%%\addbibresource{model_fidelity.bib}
%
%\renewcommand*{\bibfont}{\small}
%\newcommand{\un}[1]{\, \text{#1}}
%\newcommand{\ve}[1]{\mathbf{#1}}
%\newcommand{\veG}[1]{\text{\boldmath{$#1$}}}
%
%\begin{document}
%
%\title{Trajektorienplanung als dynamisches Optimalsteuerungsproblem und Interpretation des Lösungsraums für Fahrkomfort im automatisierten Fahren \vspace{0.2cm}}
%\subtitle{Aufgabenstellung zur Masterthesis}
%\titleinfo{Betreuer: Alexander Steinke, M.Sc.\\ \today}
%\addTitleBox{\includegraphics[width=\linewidth]{rtm_mit_schrift}}
%
%\maketitle
%%\section*{Aufgabenbeschreibung}
%
%%Zum numerischen Lösen Dynamischer Optimierungsprobleme können direkte und indirekt
%
%%Ein häufiger Ansatz zum Lösen dynamische Optimierungsprobleme ist die 
%
%%Dynamische Optimierungsprobleme lassen sich analytisch (Dynamische Optimierung) oder numerisch (Statische Optimierung) lösen. Während in der statischen Optimierung

Die allgemeine Struktur eines Optimalsteuerungsproblems lautet
%
\begin{equation*} \begin{aligned}
		\min\limits_{\ve{u}(t)} \ \ \ & J(\ve{u})=V(\ve{x}(t_\text{f}),t_\text{f}) + \int_{t_0}^{t_\text{f}}l(\ve{x}(t),\ve{u}(t),t)\text{d}t \\
		\text{u.B.v.} \ \ \ & \dot{\ve{x}}(t)=\ve{f}(\ve{x}(t),\ve{u}(t),t), \ \ve{x}(t_0)=\ve{x}_0\\ 
		& \mathbf{g}(\ve{x}(t_\text{f}),t_\text{f}) = \ve{0} \\
		& \mathbf{h}(\ve{x}(t),\ve{u}(t),t) \leq \ve{0}.
\end{aligned} \end{equation*}
%
Häufig werden dynamische Optimierungsprobleme (OP) mithilfe direkter Verfahren numerisch gelöst. 
In diesem Fall wird das dynamische in ein statisches OP überführt und der endliche Lösungsvektor $\ve{u}_\text{opt} \in \mathcal{U} \subseteq \mathbb{R}^m$ berechnet. 
Direkte Verfahren haben den Vorteil, dass Zustandsbeschränkungen leichter berücksichtigt werden können und der Konvergenzbereich größer ist. 
Indirekte Verfahren hingegen liefern eine Einsicht in die Struktur der optimalen Lösung.

In der dynamischen Optimierung hingegen werden Funktionen $\ve{u}(t)$ einer unabhängigen Variable $t$ gesucht. 
Mithilfe der Variationsrechnung können Optimalitätsbedingungen hergeleitet werden, die ein Randwertproblem formulieren. 
Die Lösung dieses Randwertproblems liefert in der Folge die optimale Steuertrajektorie $\ve{u}_\text{opt}(t)$. 
Jedoch ist das Lösen des Randwertproblems für nichtlineare Systeme häufig schwierig, weswegen auf numerische Verfahren zurückgegriffen werden muss.

Um im automatisierten Fahren (AD) gezielt Fahrprofile mittels einer optimalen Trajektorienplanung erstellen zu können, ist ein tieferes Verständnis der optimalen Lösung erforderlich. 
Mit einem direkten Verfahren ist die Interpretation der Lösung schwierig, da die Lösung aus reinen Zahlenwerten besteht. 
Zwar könnte ein gewünschtes Fahrverhalten in definierten Fahraufgaben durch zusätzliche Terme im Gütemaß und Anpassung der Gewichte erzielt werden, jedoch ist eine Übertragung auf andere Fahraufgaben fraglich. 
Eine parametrierte Steuertrajektorie $\ve{u}_\text{opt}(t)$ würde die Interpretation erheblich vereinfachen.


Ziele dieser Arbeit sind folgende Punkte:
\begin{enumerate}
	\item Übersicht für Lösungsverfahren dynamischer OPs und Einordnung des AD-Planungsproblems
	\item Erstellen von dynamischen OPs, die durch Variationsrechnung lösbar sind und Lösen dieser
	\begin{enumerate}
		\item Geeignete Fahrszenarien (z.B. Geradeausfahrt, Kurve, Geraden-Kurven-Kombination, Spurwechsel)
		\item Geeignete Komfortmerkmale (z.B. Beschleunigung, Ruck, Fahrdauer)
		\item Berücksichtigung von Begrenzungen
	\end{enumerate}
	\item Interpretation der Lösungenstrajektorien hinsichtlich des dargestellten Funktionenraums 
	\item Einordnung im Kontext Fahrkomfort
\end{enumerate}

Die Ergebnisse sind geeignet zu visualisieren und zu dokumentieren. Die aktuelle Fassung der Richtlinien zur Anfertigung von Abschlussarbeiten ist zu beachten.


%%Technische Universität Darmstadt \\
%%Institut für Automatisierungstechnik und Mechatronik \\
%%Fachgebiet Regelungstechnik und Mechatronik \\
%%Prof. Dr.-Ing. Ulrich Konigorski
%
%%\vspace*{\fill} {\tiny
%%\printbibliography[heading=none]}
%
%\end{document}


\vspace{0.5cm}
\begin{tabular}{ll}
	Beginn: 	& ------ \\
	Ende: 		& ------ \\
	Seminar: 	& ------ \\
\end{tabular}
\clearpage

\affidavit

\section*{Kurzfassung}
Zusammenfassung entsprechend der Dokumentensprache. In diesem Fall Deutsch.
\section*{Abstract}
Additional abstract in English.	

\cleardoublepage

\tableofcontents
\cleardoublepage

\mainmatter


% =================================================================================
% Kapitel und Sections etc.
% =================================================================================
\chapter{Konventionen und Beispiele}
Es folgen ein paar nützliche Beispiele und Konventionen.

\section{Gleichungen} \label{sec:gleichungen}
Bei umfangreichen Dokumenten sollte eine Gleichung, wie
%
\begin{equation} \label{equ:systemmatrix}
	\mathbf{A}_\mathrm{s} = \begin{bmatrix}1 & 0 \\ 0 & 2\end{bmatrix}
\end{equation}
%
nur eine Nummer bekommen, wenn sie, wie Gleichung~\eqref{equ:systemmatrix}, im Text referenziert wird. Die Nummerierung von Umgebungen kann über den * gesteuert werden.
% Absätze vor und nach Gleichungen vermeiden, um das Dokuemnt nicht unnötig auseinanderzuzerren

Gleichungen sind Teil des Satzes und bekommen deshalb auch Satzzeichen, wie beispielsweise die wahllosen Werte
%
\begin{align}
	\beta &= 5^{\circ}, \nonumber  \\
	c &= 42 \label{equ:antwort} 
\end{align}
%
und
\begin{equation*}
	\dot{\hat{x}} = 7 \, \mathrm{\frac{m}{s}}.
\end{equation*}
Wie an Gleichung~\eqref{equ:antwort} zu sehen, können mehrere Gleichungen im Block ausgerichtet und einzeln nummeriert und referenziert werden.

\section{Bilder und Tabellen}
Außer bei Fotos sollte auf Pixelgrafiken verzichtet werden. Vektorgrafiken können beispielsweise als PDF eingebunden werden, wie Abbildung~\ref{fig:rtm}, oder direkt mit TikZ erstellt werden.

\begin{figure} %[h] die Figure-Umgebung bietet eine Option zur Steuerung der Platzierung
	\centering
	\includegraphics[width=0.4\textwidth]{rtm_mit_schrift.pdf}
	\caption{Regelungstechnik und Mechatronik}
	\label{fig:rtm}
\end{figure}

Wie in Tabelle~\ref{tab:tabelle} zu sehen, bekommen Tabellen eine Überschrift.
%
\begin{table}
	\centering
	\caption{Tabellenüberschrift}	
	\begin{tabular}{r|cc}
		& A & B \\
		\hline
		C & 1 & 2 \\
		D & 3 & 4 \\
	\end{tabular}
	\label{tab:tabelle}
\end{table}
%

\section{Blockschaltbilder}
Blockschaltbilder können mit der im Hauptdokument eingebundenen TikZ-Bibliothek erstellt werden.
%
\begin{center} % wird von LaTex nicht frei am Seitenanfang/Ende platziert. Im Normalfall aber besser figure nutzen.
%	\centering
	\begin{tikzpicture}[node distance=5mm and 10mm]
	\node[terminal]	(w)	{};
	\node[sum]			(s) [right=of w]	{};
	\node[block]		(Regler) 		[right=of s] 			 	{Regler};
	\node[block]		(Aktor)	 		[right=of Regler]  	{Aktor};
	\node[block]		(Strecke)		[right=of Aktor]	 	{Strecke};
	\node[branch]		(punkt)			[right=of Strecke] 	{};
	\node[terminal]	(y)					[right=of punkt]		{};
	\node[block]		(Messglied) [below=of Aktor]		{Messglied};
	
	
	\draw[to] (w) -- (s) node[near start,above] {$w$};
	\draw[to] (s) -- (Regler) node[midway,above] {$e$}; 
	\draw[to] (Regler) -- (Aktor);
	\draw[to] (Aktor) -- (Strecke);
	\draw[to] (Strecke) -- (y) node[pos=0.8,above] {$y$};
	\draw[to] (punkt) |- (Messglied);
	\draw[to] (Messglied) -| (s) node[pos=0.95, right] {-};
\end{tikzpicture}
	\captionof{figure}{TikZ Blockschaltbild in Figure-Umgebung}
	\label{fig:bsb_beispiel}
\end{center}

Variablen, wie $w$ aus Abbildung~\ref{fig:bsb_beispiel}, kommen auch im Fließtext immer in eine Mathe-Umgebung.

\section{Plots}
Plots können direkt in \LaTeX mit dem Paket \texttt{pgfplots} sogar mit der TU-Farbpalette gesetzt werden.
In Abbildung~\ref{fig:step_pt2} wird z.B. die rtm-Farbe \textcolor{TUDa-2c}{\texttt{TUDa-2c}} genutzt.

Aus Matlab heraus können Plots z.B. mit dem Tool \texttt{matlab2tikz} für \LaTeX{} exportiert werden.
Der mitgelieferte \texttt{cleanfigure()}-Befehl räumt Plots vor dem Export auf und sorgt für eine kleine Dateigröße.
Mit \texttt{pgfplots} kann aber auch direkt in \LaTeX analytische Funktionen plotten: 
\tikz{	
	\begin{axis}[hide axis, height=18mm, width=5cm]
		\addplot[domain=-2*pi:2*pi, samples=100]{cos(deg(x))};
	\end{axis}
}.

\begin{figure}
	\centering
	% !TeX root = ../TUDaThesis.tex

% This file was created by matlab2tikz.
%
%The latest updates can be retrieved from
%  http://www.mathworks.com/matlabcentral/fileexchange/22022-matlab2tikz-matlab2tikz
%where you can also make suggestions and rate matlab2tikz.
%
\definecolor{mycolor1}{rgb}{0.00000,0.44700,0.74100}%
%
\begin{tikzpicture}

\begin{axis}[%
width=0.7\textwidth,
height=6cm,
scale only axis,
separate axis lines,
every outer x axis line/.append style={white!40!black},
every x tick label/.append style={font=\color{white!40!black}},
every x tick/.append style={white!40!black},
xmin=0,
xmax=14,
every outer y axis line/.append style={white!40!black},
every y tick label/.append style={font=\color{white!40!black}},
every y tick/.append style={white!40!black},
ymin=0,
ymax=2.5,
ylabel={Ausgang normiert},
xlabel={Zeit in Sekunden},
axis background/.style={fill=white}
]
\addplot [color=TUDa-2c, line width=1, forget plot]
  table[row sep=crcr]{%
0	0\\
0.0921034037197206	0.00411279805587128\\
0.184206807439441	0.0159482926069074\\
0.276310211159164	0.0347790289943219\\
0.368413614878884	0.0599132256263424\\
0.460517018598605	0.0906944372784881\\
0.552620422318325	0.126501104885605\\
0.644723826038048	0.166746003636801\\
0.736827229757768	0.210875600557069\\
0.828930633477489	0.258369332141982\\
0.921034037197209	0.308738812005936\\
1.01313744091693	0.361526977911547\\
1.10524084463665	0.416307186968996\\
1.19734424835637	0.472682267230876\\
1.38155105579581	0.588769773524355\\
1.93417147811414	0.942361883863663\\
2.11837828555358	1.05580915778143\\
2.2104816892733	1.11091852101091\\
2.30258509299302	1.16478018774458\\
2.39468849671275	1.21728576867211\\
2.48679190043247	1.26834190455422\\
2.57889530415219	1.31786932190199\\
2.67099870787191	1.36580191469028\\
2.76310211159163	1.41208585353384\\
2.85520551531135	1.45667872351849\\
2.94730891903107	1.49954869165797\\
3.03941232275079	1.54067370474229\\
3.13151572647051	1.5800407181526\\
3.22361913019023	1.61764495604156\\
3.31572253390996	1.65348920311623\\
3.40782593762968	1.68758312811146\\
3.4999293413494	1.71994263890565\\
3.59203274506912	1.75058926910707\\
3.68413614878884	1.77954959582612\\
3.77623955250856	1.80685468824809\\
3.86834295622828	1.83253958652982\\
3.960446359948	1.85664281046295\\
4.05254976366772	1.87920589727478\\
4.14465316738744	1.90027296787515\\
4.23675657110716	1.91989032080341\\
4.32885997482689	1.93810605308292\\
4.42096337854661	1.95496970715117\\
4.51306678226633	1.97053194300157\\
4.60517018598605	1.98484423464626\\
4.69727358970577	1.99795858998971\\
4.78937699342549	2.00992729318779\\
4.88148039714521	2.02080266855745\\
4.97358380086493	2.03063686509711\\
5.15779060830437	2.04738828496254\\
5.34199741574382	2.06058825850954\\
5.52620422318326	2.07063002123016\\
5.7104110306227	2.07788951993463\\
5.89461783806214	2.08272250421635\\
6.07882464550158	2.08546241975432\\
6.26303145294103	2.08641900115247\\
6.44723826038047	2.08587746840164\\
6.72354847153963	2.08281911140062\\
6.99985868269879	2.07774554823536\\
7.36827229757768	2.06896203894915\\
8.19720293105517	2.04625070133987\\
8.74982335337349	2.03209506210472\\
9.2103403719721	2.02202542105953\\
9.6708573905707	2.0138167145694\\
10.1313744091693	2.00747359450982\\
10.5918914277679	2.00284185840762\\
11.1445118500862	1.99919994883539\\
11.7892356761243	1.99698913404195\\
12.6181663096018	1.99626757654961\\
13.9997173653976	1.99742653315299\\
14.0918207691173	1.99753717612361\\
};
\addplot [color=black, dotted, forget plot]
  table[row sep=crcr]{%
-1	2\\
15.4	2\\
};
\end{axis}

\end{tikzpicture}%
	\caption{Sprungantwort $\mathrm{PT}_2$}
	\label{fig:step_pt2}
\end{figure}


% =================================================================================
% Anhang
% =================================================================================
\appendix % Damit wird der Anhang begonnen. Die Kapitel werden ab jetzt mit Buchstaben nummeriert


% =================================================================================
% Literaturverzeichnis
% =================================================================================
\cleardoublepage        % Auf eine leere Seite einfügen
\phantomsection         % Für Aufnahme ins Inhaltsverzeichnis
\addcontentsline{toc}{chapter}{\bibname}  % In Inhaltsverzeichnis von
                                          % Dokument und pdf aufnehmen
\printbibliography 
% =================================================================================

	
\end{document}
